\documentclass[11pt,a4paper]{article}
\usepackage[margin=2cm]{geometry}
\usepackage{enumitem}
\usepackage{titlesec}
\usepackage{xcolor}
\usepackage{hyperref}
\usepackage{booktabs}
\usepackage{parskip}
\usepackage{tcolorbox}

\definecolor{accent}{RGB}{0,100,150}
\definecolor{notebg}{RGB}{245,245,245}
\definecolor{noteborder}{RGB}{200,200,200}
\hypersetup{colorlinks=true, linkcolor=accent, urlcolor=accent}

\titleformat{\section}{\large\bfseries\color{accent}}{}{0em}{}[\vspace{-0.5em}\rule{\textwidth}{0.4pt}]
\titleformat{\subsection}{\normalsize\bfseries}{}{0em}{}

\setlist[enumerate]{leftmargin=*, itemsep=0.6em}

\newcommand{\notebox}{%
\begin{tcolorbox}[colback=notebg, colframe=noteborder, boxrule=0.5pt, arc=2pt, left=4pt, right=4pt, top=2pt, bottom=2pt]
\vspace{2.2cm}
\end{tcolorbox}
}

\newcommand{\bignote}[1][3.5cm]{%
\begin{tcolorbox}[colback=notebg, colframe=noteborder, boxrule=0.5pt, arc=2pt, left=4pt, right=4pt, top=2pt, bottom=2pt]
\vspace{#1}
\end{tcolorbox}
}

\newcommand{\filledbox}[1]{%
\begin{tcolorbox}[colback=notebg, colframe=noteborder, boxrule=0.5pt, arc=2pt, left=4pt, right=4pt, top=2pt, bottom=2pt]
#1
\end{tcolorbox}
}

\begin{document}

\begin{center}
{\LARGE\bfseries Voorbereiding gesprek}\\[0.3em]
{\large WAVES collega, IOF \& VLAIO Innovatiemandaat}\\[0.3em]
{\small 2 april 2026}
\end{center}

\vspace{0.5em}

\begin{tcolorbox}[colback=white, colframe=accent, boxrule=1pt, arc=3pt, title={\textbf{Algemene notities}}, fonttitle=\bfseries\color{white}, colbacktitle=accent]
\begin{itemize}[leftmargin=*, itemsep=0.3em]
\item IOF deadlines 2026: 21/4 (maart-call), 3/8 (juni-call), 19/10 (september-call).
\item WAVES is UGent/imec maar IP is eigenlijk van UGent. imec geeft ``een gunst'' van het IP te hebben.
\item Luc heeft indertijd ook een spin-off gedaan: Excentis.
\item KBC Start It is ook interessant: kantoren gebruiken, minder funding, meer network/mentorship. Bespreken met Gentrepreneur of online zoeken.
\item ``Ik heb al veel onderzoeksvragen gedaan'' (ze had al veel research achter de rug voor IOF).
\end{itemize}
\end{tcolorbox}

%----------------------------------------------------------------------
\section{Belangrijk}
%----------------------------------------------------------------------

\subsection{De grote vragen}

\begin{enumerate}

\item \textbf{Heb jij een co-founder?} Hoe heb je dat aangepakt? istart vereist minstens 2 founders voor het hoofdprogramma. Heb je iemand specifiek gezocht of had je al iemand? Hoe heb je de equity verdeeld?

\filledbox{Nog geen co-founder. Ze zoekt nu een Entrepreneur in Residence (EiR) die meer de business kant doet.}

\item \textbf{Ik overweeg mijn vriendin als co-founder.} Ze zit in dezelfde onderzoeksgroep, zelfde promotor, complementair domein (EM effecten op neuraal weefsel). Maar het blijft een koppel dat samen een bedrijf start. Ken je voorbeelden hiervan? Is dat een probleem bij istart of VLAIO? Hoe kijkt een jury daarnaar?

\filledbox{imec istart mensen moeten echt involved zijn, dus het idee met vriendin is ``op het randje.'' Niet per se afgekeurd, maar ze moeten echt betrokken zijn bij het bedrijf.}

\item \textbf{Hoe combineer je een PhD met een spin-off voorbereiden?} Ik ben bijna klaar met mijn doctoraat, maar Carolina zit er nog middenin. Kan zij co-founder zijn terwijl ze haar PhD doet? Zijn daar regels over bij UGent?

\notebox

\item \textbf{Eerlijk: is het het waard?} Terugkijkend, ben je blij dat je het spin-off pad gekozen hebt? Zijn er momenten geweest dat je dacht ``ik had beter gewoon in de industrie gegaan''?

\notebox

\item \textbf{Hoeveel werk is het echt?} Ik hoor dat het mandaat vrij flexibel is qua uren. Hoeveel uur per week ben je er effectief mee bezig? Kan je nog andere dingen doen ernaast?

\notebox

\item \textbf{De industriele mentor: moet dat echt een bestaand bedrijf zijn?} Ik heb geen band met een bedrijf op dit moment. Luc legt contact met Ericsson, maar dat is pril. Hoe erg is het als je mentor niet uit je exacte domein komt?

\filledbox{Ze koos haar hardware partner als mentor, iemand betrokken bij de ontwikkeling waar ze het goed mee overeen kwam. Eerlijk gezegd was het ``voor de vorm,'' niet zo belangrijk. Maar de persoon moet wel echt bij een bedrijf zijn met echte bedrijfservaring.}

\end{enumerate}

\subsection{De VLAIO-juryverdediging}

\begin{enumerate}

\item \textbf{Hoe was de jury samengesteld?} VLAIO zegt ``externe experten'' die op twee assen beoordelen: wetenschappelijke kwaliteit en valorisatie. Waren dat aparte mensen (academici vs.\ industrie)? Hoeveel juryleden? Zat er iemand van VLAIO zelf bij?

De scores gaan blijkbaar van kritisch / zwak / redelijk / goed / uitmuntend. Het maakt nogal uit wie die scores geeft.

\filledbox{Bij IOF: externe experten uit het domein (in haar geval dierenartsen). Ook iemand die zelf een business in de paardenwereld had. Mix van vaste leden en externe experten.}

\item \textbf{Wat voor vragen stelden ze?} Vooral technisch (``bewijs dat het werkt'') of vooral commercieel (``wie koopt dit'')? Welke as kreeg meer pushback?

\filledbox{Niet enkel valorizatie, ook technische vragen. Maar valorizatiestrategie is het belangrijkste punt bij IOF: je moet laten zien hoeveel geld en hoeveel tijd ervoor nodig is.}

\item \textbf{Was er een vraag die je niet had verwacht?}

\notebox

\item \textbf{Hoe lang was de presentatie vs.\ de vragenronde?} Totale tijd en verdeling? Ik hoorde dat het zoiets als 10 minuten is.

\filledbox{Heeft al een template voor de slides.}

\item \textbf{Kan ik je slides zien?} Wout zei dat ik je ernaar moest vragen. Ik wil graag zien hoe je het hebt opgebouwd, hoeveel slides, hoeveel tekst vs.\ figuren, hoe technisch vs.\ commercieel.

\filledbox{Heeft een template. Opvolgen: vragen om die door te sturen.}

\end{enumerate}

\subsection{Slaagkans en competitiviteit}

\begin{enumerate}[resume]

\item \textbf{Hoe competitief was jouw call?} De tweede call van 2025 had 25 aanvragen en 11 goedkeuringen, dus ongeveer 44\%. Was dat bij jou vergelijkbaar?

VLAIO publiceert geen historische slaagcijfers, dus dat is het enige publieke datapunt dat ik heb.

\notebox

\item \textbf{Was het budget de beperking, of de kwaliteit?} Werden goede projecten afgewezen omdat het geld op was, of werd alles boven een kwaliteitsdrempel gefinancierd?

\notebox

\end{enumerate}

\subsection{UGent TechTransfer en IP}

\begin{enumerate}[resume]

\item \textbf{Wat neemt UGent precies bij een spin-off?} De Fast Lane (goedgekeurd maart 2025) belooft ``vooraf bepaalde standaardvoorwaarden,'' maar de concrete cijfers staan nergens. Welk equity-percentage? Royalty? Success fee?

Er zijn blijkbaar twee tracks: een gewone met equity, en een ``Mini'' met enkel royalty + success fee. Ik heb de echte cijfers nodig.

\filledbox{\textbf{Fast Lane cijfers: 0.5\% royalty en 6\% equity.} Ze raadt Fast Lane aan omdat TechTransfer super traag is met onderhandelingen. imec istart vindt dat niet fijn en wil liever niet wachten op TT. Alternatief: spin-off aanvragen bij TT voor 6\% en dan laten ze je met rust. OF IOF aanvragen, ook 6\%, met salaris dat eigenlijk postdoc barema is.}

\item \textbf{Hoe lang duurde de IP/licentie-onderhandeling?} Heb je de Fast Lane gebruikt of het oude model? Weken? Maanden?

\filledbox{TechTransfer is super traag met onderhandelingen. Geen concreet tijdsbestek genoemd, maar traagheid was een reden om Fast Lane te verkiezen.}

\item \textbf{Iets in het TechTransfer-proces dat je verraste of dat je graag op voorhand had geweten?}

\filledbox{TechTransfer zijn meer dan juristen. Ze hebben haar teksten verbeterd (de geschreven aanvraag). Positieve verrassing.}

\end{enumerate}

\subsection{De industriele mentor}

\begin{enumerate}[resume]

\item \textbf{Hoe heb je je industriele mentor gevonden?} Kende je die al, of heeft TechTransfer of VLAIO jullie gekoppeld?

VLAIO eist een getekende samenwerkingsovereenkomst tussen UGent en het bedrijf van de mentor binnen 4 weken na indiening. Bij spin-off mandaten mag de mentor ook een VC-fonds of consultancybureau zijn. Ik moet er voor september eentje vinden.

\filledbox{Ze koos haar hardware partner. Was al betrokken bij de ontwikkeling en ze kwam er goed mee overeen. Kende die dus al zelf.}

\item \textbf{Hoe betrokken is de mentor in de praktijk?} Maandelijks bellen? Per kwartaal? Of eerder een handtekening op de rapporten?

\filledbox{Eerlijk gezegd ``voor de vorm.'' Niet zo belangrijk in de praktijk.}

\item \textbf{Vroeg de jury naar de mentor tijdens de verdediging?} Maakte het profiel van de mentor (senioriteit, bedrijf, domein) uit voor de beoordeling?

\notebox

\end{enumerate}

%----------------------------------------------------------------------
\section{Minder belangrijk}
%----------------------------------------------------------------------

\subsection{De mandaatperiode in de praktijk}

\begin{enumerate}[resume]

\item \textbf{Hoeveel vrijheid had je in hoe je je tijd besteedde?} VLAIO zegt ``meerderheid van de tijd voor onderzoek.'' Maar op de spin-off track zijn klantgesprekken en business development het hele punt. Hoe ben je daarmee omgegaan? Controleert iemand dat?

\filledbox{IOF is echt om onderzoeksvragen te beantwoorden, maar niets tegen om ook business onderzoek te doen. WPs kunnen ook valorizatie zijn, bvb business model toetsen aan de realiteit. Ze mag openlijk commerciele activiteiten doen.}

\item \textbf{Hoe zijn de halfjaarlijkse voortgangsrapporten?} Hoe gedetailleerd? Hoe lang? Wie leest ze? Stressvol of formaliteit?

\notebox

\item \textbf{Kon je openlijk commerciele activiteiten doen tijdens het mandaat?} Klantbezoeken, pitchen op conferenties, gesprekken met investeerders?

\filledbox{Ja, mag openlijk.}

\item \textbf{Wanneer heb je het bedrijf opgericht?} Het mandaat stopt formeel bij oprichting van de spin-off, dus er is een prikkel om uit te stellen. Maar istart wil misschien een bestaand bedrijf. Wat is de juiste timing?

\filledbox{Bij imec istart moet je meteen incorporeren (bedrijf starten). Ze krijg je een lump sum + equity en dan moet je meteen het bedrijf oprichten.}

\item \textbf{Kan je het VLAIO-mandaat en imec.istart tegelijk doen?} Zijn ze compatibel of sluit het ene het andere uit? Heeft iemand in jouw cohort ze gecombineerd?

\filledbox{IOF kun je met VLAIO stapelen, maar NIET met imec istart. istart wil dat je meteen het bedrijf start. Ze is niet zeker of VLAIO daar een probleem mee heeft, maar IOF wel (conflict met istart).}

\end{enumerate}

\subsection{Het IOF-traject}

\begin{enumerate}[resume]

\item \textbf{Welk IOF-traject heb je gebruikt?} ConcepTT / StarTT / Advanced / Stepstone? Waarom dat?

Deadlines 2026: 21/4 (maart-call), 3/8 (juni-call), 19/10 (september-call). Ik denk aan StarTT of ConcepTT om te beginnen, Stepstone later als ik mensen moet aanwerven.

\filledbox{Ze startte eerst met IOF (vermoedelijk StarTT/ConcepTT) en deed daarna een Advanced. Het traject ``wordt eigenlijk opgelegd'' (de IOF-commissie adviseert welk traject past).}

\item \textbf{Hoe was de intake meeting?} Wie zat erbij? Formaliteit of echte evaluatie?

\filledbox{Intake meeting is super belangrijk! Goed voorbereiden. Valorizatie moet al gedaan zijn op dat moment.}

\item \textbf{Hoeveel IOF-geld heb je gekregen, en welke voorwaarden?} Rapportering? Bestedingsbeperkingen? Tijdslijn?

\filledbox{``Zij was het duurste.'' Heeft er experimenten mee gedaan. Geen concreet bedrag genoemd.}

\item \textbf{Kan je IOF en VLAIO tegelijk lopen?} Overlap of conflict?

\filledbox{Ja, IOF en VLAIO kun je stapelen. Geen conflict.}

\item \textbf{Was de IOF-aanvraag makkelijker of moeilijker dan VLAIO?} Andere evaluatiestijl?

\notebox

\end{enumerate}

\subsection{imec.istart}

\begin{enumerate}[resume]

\item \textbf{Weet je iets over imec.istart vanuit collega's of eigen ervaring?} Iemand in WAVES die het heeft gedaan?

\filledbox{Ze wou zelf niet imec istart doen. Redenen: (1) je geeft equity op voor funding die minder is dan wat UGent biedt (``UGent is prima''), (2) ze moest dan een co-founder vinden, (3) ze had nog onderzoeksvragen. Belangrijk: de istart equity is niet dilutable!

imec istart is 100K pre-seed: 50K voor 6\% equity + 50K convertible loan. Ze willen een PoC. Eerst binnengeraken, dan voorbereiding op pitching, en dan moet je nog naar investeerders omdat 100K niet genoeg is. Je krijgt het imec network, dat verantwoordt het wel.}

\item \textbf{De Media, Telecom \& Entertainment vertical:} draait die nog? Een ouder istart-handboek vermeldt Telenet als partner. Weet je of dat nog klopt?

\notebox

\item \textbf{De Expert-in-Residence:} weet je hoe dat in de praktijk werkt? Is het echt 1 a 2 dagen per week ingebed, of meer lichte begeleiding?

\filledbox{Ze zoekt nu zelf een Entrepreneur in Residence (EiR). Die zou meer de business kant doen. Geen details over de dagelijkse werking.}

\item \textbf{Enig idee wat het istart Investment Committee anders zoekt dan een VLAIO-panel?} Meer commercieel? Meer technisch?

\notebox

\end{enumerate}

\subsection{Gentrepreneur en IOF-details}

\begin{enumerate}[resume]

\item \textbf{Kan je IOF Stepstone en het IOF Spin-off traject stapelen?} De brochure toont Stepstone (tot 400K) en een apart Spin-off traject (50 tot 250K). Zijn die sequentieel, of kan je beide krijgen?

\filledbox{Niet direct beantwoord. Wel: je kunt een spin-off aanvragen bij TT voor 6\% en dan laten ze je met rust, OF je kunt IOF aanvragen (ook 6\%) met salaris dat eigenlijk postdoc barema is.}

\item \textbf{Heb je de Gentrepreneur Research Track gedaan?} Was dat nuttig voor iemand die al een spin-off plant, of te basaal?

\filledbox{Echt goed! 4 dagen. Je leert je business plan schrijven. Je bespreekt veel nuttige dingen: IP, business plan, goede eerste inleiding. Redelijk lightweight. Aanrader.}

\item \textbf{De Spin-Off Caf\'e:} ben je er geweest? Echt nuttig om mentoren of cofounders te vinden, of eerder een netwerkdrink?

\filledbox{Super waardevol. Meer perspectieven. Veel ondersteuning.}

\end{enumerate}

%----------------------------------------------------------------------
\section{De rest}
%----------------------------------------------------------------------

\subsection{Politiek en praktisch}

\begin{enumerate}[resume]

\item \textbf{Hoe stond Wout tegenover het spin-off traject?} Spanning tussen ``onderzoek afmaken'' en ``bedrijf bouwen''?

\notebox

\item \textbf{Is er iemand bij VLAIO die je informeel kan bellen voor je indient?} Om de aanvraag te checken of te peilen wat de jury zoekt?

Ik ken de UGent-contacten (Fien Elderweirdt, Nina De Wilde, Karen Cur\'e bij TechTransfer). Maar is er iemand aan de VLAIO-kant die bereikbaar is?

\notebox

\item \textbf{Iemand in WAVES of TechTransfer die bijzonder nuttig was?} Een specifiek persoon waar ik vroeg naartoe moet.

\notebox

\item \textbf{De infosessie in januari} waar VLAIO-coordinatoren de call toelichten aan UGent: de moeite of overslaan?

\notebox

\item \textbf{Terugkijkend: wat zou je anders doen?}

\bignote

\item \textbf{En: wat ging beter dan verwacht?}

\bignote

\end{enumerate}

\subsection{Financieel en contractueel}

\begin{enumerate}[resume]

\item \textbf{Wat was je effectieve loon op het mandaat?} Standaard UGent postdoc barema? Welk?

\filledbox{Eigenlijk postdoc barema (bij IOF). Geen specifiek schaal/trap genoemd.}

\item \textbf{Heeft VLAIO ooit iets teruggeduwd tijdens het mandaat?} Een budgetpost, een koerswijziging, een voortgangsrapport?

\notebox

\item \textbf{Verborgen kosten?} Octrooikosten, juridische kosten voor de samenwerkingsovereenkomst, iets uit eigen zak?

\filledbox{Patenten: universiteit heeft ze graag, maar het kost veel om te onderhouden. Een patent is enkel voor Belgie, je moet het globaal/EU doen om echt te verdedigen. En je beschrijft alles, dus alles staat er in. Als een ander bedrijf je kopieert moet je die ook aanklagen. Het is haar afgeraden omdat het veel geld kost en je moet alles vrijgeven. Maar geeft wel credibiliteit voor investeerders. Ze werkt nu met trade secrets. Klanten hebben wel graag dat het wetenschappelijk onderbouwd is.}

\end{enumerate}

\subsection{De geschreven aanvraag}

\begin{enumerate}[resume]

\item \textbf{Hoe lang was de geschreven aanvraag?} Is er een paginalimiet?

\notebox

\item \textbf{Wat was de balans tussen onderzoeksplan en businessplan?} Welk deel woog zwaarder in de evaluatie?

\filledbox{Bij IOF is valorizatie heel belangrijk. Eigenlijk voor ALLE trajecten (VLAIO, imec istart): hoe ga je geld verdienen, hoeveel, wie is mijn klant, mijn stakeholders, kosten product. Het belangrijkste punt is de valorizatiestrategie.}

\item \textbf{Had je klantbewijs bijgevoegd?} LOIs, gesprekken, marktonderzoek? Vond de jury dat belangrijk?

\filledbox{Dit is moeilijk, maar je moet echt al gaan spreken met mensen. Heel duidelijk beeld: hoeveel is de markt, hoeveel mensen willen ervoor betalen. Concurrenten geven niet graag cijfers vrij. Volg de competitie op LinkedIn, zoek hun cijfers. Zo heeft ze eens bij een concurrent gevonden hoeveel euro per maand en hoeveel klanten in zoveel maanden tijd.}

\item \textbf{Wie heeft je geholpen met schrijven?} TechTransfer, Team Vlaams-federaal, je promotor? Wie gaf de nuttigste feedback?

\filledbox{TechTransfer heeft haar teksten verbeterd. Ze zijn meer dan juristen.}

\end{enumerate}

%----------------------------------------------------------------------
\vspace{1em}
\begin{tcolorbox}[colback=white, colframe=accent, boxrule=1pt, arc=3pt, title={\textbf{Opvolging}}, fonttitle=\bfseries\color{white}, colbacktitle=accent]
\begin{itemize}[leftmargin=*, itemsep=0.3em]
\item Slides template opvragen
\item IOF deadline 21/4 checken (is dat haalbaar?)
\item Gentrepreneur Research Track inschrijven
\item Spin-Off Caf\'e bijwonen
\item KBC Start It bekijken
\item Valorizatiestrategie uitwerken voor intake meeting
\item Concurrenten opzoeken op LinkedIn (prijzen, klanten)
\end{itemize}
\end{tcolorbox}

\vspace{0.5em}
\noindent\textbf{Voor het afsluiten:}
\begin{itemize}[leftmargin=*]
\item Vragen of ze bereid is een draft VLAIO-aanvraag na te lezen voor indiening.
\item Email of nummer voor opvolgvragen.
\end{itemize}

\end{document}
